\section{Maximum Sum Combination}
\label{sec:heap-max-sum-combination}

\problemstatement{We are given two arrays $A$ and $B$, each of size $N$,
and an integer $K$. Consider every possible pair sum $A[i]+B[j]$ for
$0 \le i,j < N$ (there are $N^2$ such pairs). Find the $K$ \textbf{largest}
pair sums, in descending order, \textbf{without} computing all $N^2$ sums.}

\begin{patternbox}
\emph{``$K$ largest sums formed by combining one element from each of two
(sorted) collections''} is a generalization of Kth Largest
(Section~\ref{sec:heap-kth-largest}) to a combinatorially large candidate
space: there are $N^2$ possible pairs, far too many to generate and sort
when $K \ll N^2$. The keyword to notice is that the candidate sums have
\textbf{internal structure} (they come from a grid of sorted index pairs),
which a heap can exploit to generate only the $K$ largest sums directly,
without ever materializing the rest.
\end{patternbox}

\subsection*{Understanding the problem carefully}

Sort both $A$ and $B$ in \textbf{descending} order first. Now,
$A[0]+B[0]$ is guaranteed to be the single largest possible sum (the
largest element of each array, paired together). Crucially, once we know
$A[0]+B[0]$ is the maximum, the \emph{next}-largest sum must come from
either decreasing the $A$-index to $1$ (pairing $A[1]$ with $B[0]$) or
decreasing the $B$-index to $1$ (pairing $A[0]$ with $B[1]$) -- no other
pair can exceed both of these, since both arrays are sorted descending,
and any pair using indices further from $(0,0)$ in either coordinate can
only be smaller than at least one of these two candidates.

\begin{ideabox}[Greedy Choice, Powered by a Max-Heap over Index Pairs]
Maintain a max-heap of candidate \emph{index pairs} $(i,j)$, keyed by
$A[i]+B[j]$, along with a \textit{visited} set to avoid pushing the same
pair twice. Start with $(0,0)$ in the heap. Repeat $K$ times: pop the pair
with the largest sum, record that sum, then push $(i+1,j)$ and $(i,j+1)$
(if within bounds and not already visited) as the next candidates that
could potentially be the next-largest sum.
\end{ideabox}

\subsection*{Why only two neighbors need to be considered at each step}

This is the same idea as merging $k$ sorted lists
(Section~\ref{sec:heap-merge-k-lists}), generalized to a 2D grid of sorted
candidates rather than $k$ separate 1D lists: because both $A$ and $B$ are
sorted descending, the sum $A[i]+B[j]$ is non-increasing as $i$ or $j$
increases (holding the other fixed) -- so the grid of all possible sums is
itself ``sorted'' along both axes, exactly like the row/column-sorted
matrices in Sections~\ref{sec:bs-search-2d-matrix}--\ref{sec:bs-matrix-median}.
Once a pair $(i,j)$ has been popped as the current maximum, the only pairs
that could possibly be the \emph{next}-largest, among those not yet
considered, are its immediate neighbors $(i+1,j)$ and $(i,j+1)$ -- every
other still-uncovered pair is reachable only by passing through one of
these two first, so they are never skipped over by this expansion rule.

\subsection*{Algorithm}

\begin{enumerate}
  \item Sort $A$ and $B$ in descending order.
  \item Initialize a max-heap (by sum) with the single pair $(0,0)$;
        mark $(0,0)$ visited.
  \item Repeat $K$ times:
  \begin{enumerate}
    \item Pop $(i,j)$ with the largest sum $A[i]+B[j]$; record this sum.
    \item If $i+1 < N$ and $(i+1,j)$ not visited: push it; mark
          visited.
    \item If $j+1 < N$ and $(i,j+1)$ not visited: push it; mark
          visited.
  \end{enumerate}
  \item Return the $K$ recorded sums.
\end{enumerate}

\begin{algorithm}[H]
\caption{Maximum Sum Combination}
\begin{algorithmic}[1]
\Procedure{MaxSumCombination}{$A, B, K$}
  \State sort $A, B$ descending
  \State $\textit{heap} \gets$ max-heap with $(A[0]+B[0], 0, 0)$
  \State mark $(0,0)$ visited
  \State $\textit{result} \gets$ empty list
  \For{$\textit{count} \gets 1$ to $K$}
    \State $(\textit{sum}, i, j) \gets$ pop maximum
    \State append \textit{sum} to \textit{result}
    \If{$i+1<N$ \textbf{and} $(i{+}1,j)$ not visited}
      \State push $(A[i{+}1]+B[j], i{+}1, j)$; mark visited
    \EndIf
    \If{$j+1<N$ \textbf{and} $(i,j{+}1)$ not visited}
      \State push $(A[i]+B[j{+}1], i, j{+}1)$; mark visited
    \EndIf
  \EndFor
  \State \Return \textit{result}
\EndProcedure
\end{algorithmic}
\end{algorithm}

\subsection*{Dry run}

\dryrun{Take $A = [3,2]$, $B = [1,4]$, $K = 2$. Sorted descending:
$A=[3,2]$, $B=[4,1]$.}

\begin{itemize}
  \item Start: heap $=\{(7,0,0)\}$ (sum $A[0]+B[0]=3+4=7$). Visited
        $=\{(0,0)\}$.
  \item Pop $(7,0,0)$: record $7$. Push $(1,0)$: sum
        $A[1]+B[0]=2+4=6$. Push $(0,1)$: sum $A[0]+B[1]=3+1=4$. Heap
        $=\{6,4\}$. Visited $=\{(0,0),(1,0),(0,1)\}$.
  \item Pop $(6,1,0)$: record $6$. $K=2$ reached; stop.
\end{itemize}

The top $2$ sums are $[7,6]$. Checking by hand: all four pair sums are
$3{+}4=7$, $3{+}1=4$, $2{+}4=6$, $2{+}1=3$; sorted descending,
$[7,6,4,3]$, and the top $2$ are indeed $[7,6]$.

\subsection*{C++ implementation}

\begin{lstlisting}
#include <bits/stdc++.h>
using namespace std;

vector<int> maxSumCombination(vector<int> A, vector<int> B, int K) {
    int n = A.size();
    sort(A.rbegin(), A.rend());
    sort(B.rbegin(), B.rend());

    // (sum, i, j)
    priority_queue<tuple<int,int,int>> maxHeap;
    set<pair<int,int>> visited;

    maxHeap.push({A[0] + B[0], 0, 0});
    visited.insert({0, 0});

    vector<int> result;
    while (K-- > 0 && !maxHeap.empty()) {
        auto [sum, i, j] = maxHeap.top();
        maxHeap.pop();
        result.push_back(sum);

        if (i + 1 < n && !visited.count({i+1, j})) {
            maxHeap.push({A[i+1] + B[j], i+1, j});
            visited.insert({i+1, j});
        }
        if (j + 1 < n && !visited.count({i, j+1})) {
            maxHeap.push({A[i] + B[j+1], i, j+1});
            visited.insert({i, j+1});
        }
    }
    return result;
}

int main() {
    vector<int> A = {3, 2};
    vector<int> B = {1, 4};
    for (int x : maxSumCombination(A, B, 2)) cout << x << " ";
    cout << endl;
    return 0;
}
\end{lstlisting}

\begin{complexitybox}
\textbf{Time Complexity.} Sorting costs $O(N \log N)$. The main loop runs
$K$ times, each doing $O(1)$ pushes/pops of $O(\log K)$ cost (the heap
never holds more than $O(K)$ candidate pairs), giving an overall time
complexity of
\[
  O(N \log N + K \log K).
\]
\textbf{Space Complexity.} $O(K)$ for the heap and visited set.
\end{complexitybox}

\begin{takeawaybox}
``$K$ largest/smallest results from combining elements of two or more
sorted collections'' generalizes the merge-$k$-sorted-lists idea
(Section~\ref{sec:heap-merge-k-lists}) to a combinatorial candidate space:
instead of $k$ separate sequences, the candidates form an implicit sorted
grid, and a heap explores that grid one ``neighbor expansion'' at a time,
from the best candidate outward, generating exactly the $K$ best results
without ever enumerating the rest.
\end{takeawaybox}
